% answers.tex —— 2012 年·江西卷·文科 · 参考答案与详细解析片段
% 由 tools/generate_projects.py 依来源总账自 2012/2012_数学真题_江西文_试题.md + 2012/2012_数学真题_江西文_解析正文.md + 2012/2012_数学真题_江西文_答案要点.md 生成，勿手改（改源 MD 后重跑）

\begin{answer}
\gksolhead{1}{}

\textbf{题目}：

若复数 $z=1+\mathrm i$（$\mathrm i$ 为虚数单位），$\overline z$ 是 $z$ 的共轭复数，则 $z^2+\overline z^2$ 的虚部为\mbox{（　　）}

（A）$0$；（B）$-1$；（C）$1$；（D）$-2$。

\Needspace{4\baselineskip}
\textbf{解析}
$z^2=(1+\mathrm i)^2=2\mathrm i$，$\overline z^2=(1-\mathrm i)^2=-2\mathrm i$，故 $z^2+\overline z^2=0$，虚部为 $0$。选 A。

\gkanswer{答案：A。}

\end{answer}
\begin{answer}
\gksolhead{2}{}

\textbf{题目}：

若全集 $U=\{x\in\mathbf R\mid x^2\leqslant4\}$，则集合 $A=\{x\in\mathbf R\mid|x+1|\leqslant1\}$ 的补集 $\complement_UA$ 为\mbox{（　　）}

（A）$\{x\in\mathbf R\mid0<x<2\}$；（B）$\{x\in\mathbf R\mid0\leqslant x<2\}$；（C）$\{x\in\mathbf R\mid0<x\leqslant2\}$；（D）$\{x\in\mathbf R\mid0\leqslant x\leqslant2\}$。

\Needspace{4\baselineskip}
\textbf{解析}
$U=[-2,2]$，$A=[-2,0]$，故 $\complement_UA=(0,2]$。选 C。

\gkanswer{答案：C。}

\end{answer}
\begin{answer}
\gksolhead{3}{}

\textbf{题目}：

设函数 $f(x)=\begin{cases}x^2+1,&x\leqslant1,\\\dfrac2x,&x>1,\end{cases}$ 则 $f(f(3))=$\mbox{（　　）}

（A）$\dfrac15$；（B）$3$；（C）$\dfrac23$；（D）$\dfrac{13}9$。

\Needspace{4\baselineskip}
\textbf{解析}
$f(3)=\dfrac23\leqslant1$，故 $f(f(3))=f\left(\dfrac23\right)=\dfrac49+1=\dfrac{13}9$。选 D。

\gkanswer{答案：D。}

\end{answer}
\begin{answer}
\gksolhead{4}{}

\textbf{题目}：

若 $\dfrac{\sin\alpha+\cos\alpha}{\sin\alpha-\cos\alpha}=\dfrac12$，则 $\tan2\alpha=$\mbox{（　　）}

（A）$-\dfrac34$；（B）$\dfrac34$；（C）$-\dfrac43$；（D）$\dfrac43$。

\Needspace{4\baselineskip}
\textbf{解析}
\Needspace{5\baselineskip}
由 $\dfrac{\sin\alpha+\cos\alpha}{\sin\alpha-\cos\alpha}=\dfrac12$ 得 $2\sin\alpha+2\cos\alpha=\sin\alpha-\cos\alpha$，即 $\tan\alpha=-3$。于是
\nopagebreak[4]
\[\tan2\alpha=\frac{2\tan\alpha}{1-\tan^2\alpha}=\frac{-6}{1-9}=\frac34.\]
选 B。

\gkanswer{答案：B。}

\end{answer}
\begin{answer}
\gksolhead{5}{}

\textbf{题目}：

观察下列事实：$|x|+|y|=1$ 的不同整数解 $(x,y)$ 的个数为 $4$，$|x|+|y|=2$ 的不同整数解 $(x,y)$ 的个数为 $8$，$|x|+|y|=3$ 的不同整数解 $(x,y)$ 的个数为 $12$，$\cdots$，则 $|x|+|y|=20$ 的不同整数解 $(x,y)$ 的个数为\mbox{（　　）}

（A）$76$；（B）$80$；（C）$86$；（D）$92$。

\Needspace{4\baselineskip}
\textbf{解析}
$|x|+|y|=n$ 的整数解个数为 $4n$（$n\in\mathbf N^*$），故 $n=20$ 时有 $80$ 个。选 B。

\gkanswer{答案：B。}

\end{answer}
\begin{answer}
\gksolhead{6}{}

\textbf{题目}：

小波一星期的总开支分布如图 1 所示，一星期的食品开支如图 2 所示，则小波一星期的鸡蛋开支占总开支的百分比为\mbox{（　　）}
\nopagebreak[4]
\begin{center}% 2012 年江西卷文科第 6 题：总开支饼图（图 1）与食品开支柱图（图 2）。
% 来源 PDF：普通高考/2012/2012江西文.pdf，第 1 页（印面 749）左栏。
% 图 1 标注：食品 30%、储蓄 20%、娱乐 10%、日常 10%、通讯 5%、其他 5%，
% 其余 20% 在原图中未标注，按原样留白，不补文字。
\begin{tikzpicture}[semithick, font=\small]
  % ---------- 图 1 饼图 ----------
  \begin{scope}
    \coordinate (O) at (0,0);
    \def\R{1.1}
    \fill[gray!12] (O) -- (90:\R) arc (90:-18:\R) -- cycle;
    \fill[gray!25] (O) -- (-18:\R) arc (-18:-90:\R) -- cycle;
    \fill[gray!12] (O) -- (-90:\R) arc (-90:-126:\R) -- cycle;
    \fill[gray!25] (O) -- (-126:\R) arc (-126:-162:\R) -- cycle;
    \fill[gray!12] (O) -- (-162:\R) arc (-162:-180:\R) -- cycle;
    \fill[gray!25] (O) -- (162:\R) arc (162:144:\R) -- cycle;
    \draw (O) circle (\R);
    \foreach \a in {90,-18,-90,-126,-162,-180,162,144} \draw (O) -- (\a:\R);
    \node[above right] at (36:1.62) {食品 $30\%$};
    \node[right] at (-54:1.62) {储蓄 $20\%$};
    \node[below left] at (-112:1.62) {娱乐 $10\%$};
    \node[below right] at (-150:1.62) {日常 $10\%$};
    \node[left] at (185:1.55) {通讯 $5\%$};
    \node[above left] at (150:1.55) {其他 $5\%$};
    \node[below] at (0,-2.35) {图 1};
  \end{scope}
  % ---------- 图 2 柱状图 ----------
  \begin{scope}[shift={(6.6,-0.3)}]
    \draw[->] (0,0) -- (6.2,0) node[below] {元};
    \draw[->] (0,0) -- (0,4.3);
    \foreach \y/\t in {1/40,2/80,3/120} \draw (0,\y) -- (-0.1,\y) node[left] {\t};
    \foreach \x/\h/\lab in {1/30/鸡蛋, 2/40/牛奶, 3/100/肉类, 4/80/蔬菜, 5/50/其他} {
      \draw (\x-0.3,0) rectangle (\x+0.3,\h/40);
      \node[above] at (\x,\h/40+0.12) {\lab};
    }
    \node[below] at (3,-0.8) {图 2};
  \end{scope}
\end{tikzpicture}
\end{center}
\nopagebreak[4]
（A）$30\%$；（B）$10\%$；（C）$3\%$；（D）不能确定。

\Needspace{4\baselineskip}
\textbf{解析}
\Needspace{5\baselineskip}
由图 2，食品总开支为 $30+40+100+80+50=300$ 元，其中鸡蛋 $30$ 元，占食品开支的 $\dfrac{30}{300}=10\%$。又食品开支占总开支的 $30\%$，故鸡蛋占总开支的
\nopagebreak[4]
\[10\%\times30\%=3\%.\]
选 C。

\gkanswer{答案：C。}

\end{answer}
\begin{answer}
\gksolhead{7}{}

\textbf{题目}：

若一个几何体的三视图如图所示，则此几何体的体积为\mbox{（　　）}
\nopagebreak[4]
\begin{center}% 2012 年江西卷文科第 7 题：主视图 3 个 1×1 方格、左视图 2 个、俯视图为六边形。
% 来源 PDF：普通高考/2012/2012江西文.pdf，第 1 页（印面 749）中栏。
% 六边形顶点 (1,2),(2,2),(3,1),(2,0),(1,0),(0,1)，面积 4；棱柱高 1，体积 4。
\begin{tikzpicture}[scale=0.95, semithick, font=\small]
  % 主视图
  \begin{scope}
    \draw (0,0) rectangle (3,1);
    \draw (1,0) -- (1,1); \draw (2,0) -- (2,1);
    \draw[<->] (0,-0.28) -- (1,-0.28); \node[below] at (0.5,-0.28) {1};
    \draw[<->] (1,-0.28) -- (2,-0.28); \node[below] at (1.5,-0.28) {1};
    \draw[<->] (2,-0.28) -- (3,-0.28); \node[below] at (2.5,-0.28) {1};
    \node[below] at (1.5,-0.95) {主视图};
  \end{scope}
  % 左视图
  \begin{scope}[shift={(5.2,0)}]
    \draw (0,0) rectangle (2,1);
    \draw (1,0) -- (1,1);
    \draw[<->] (0,-0.28) -- (1,-0.28); \node[below] at (0.5,-0.28) {1};
    \draw[<->] (1,-0.28) -- (2,-0.28); \node[below] at (1.5,-0.28) {1};
    \draw[<->] (2.35,0) -- (2.35,1); \node[right] at (2.35,0.5) {1};
    \node[below] at (1,-0.95) {左视图};
  \end{scope}
  % 俯视图
  \begin{scope}[shift={(0.2,-4.7)}]
    \draw (1,2) -- (2,2) -- (3,1) -- (2,0) -- (1,0) -- (0,1) -- cycle;
    \draw[<->] (1,2.3) -- (2,2.3); \node[above] at (1.5,2.3) {1};
    \draw[<->] (0.45,-0.3) -- (1,-0.3); \node[below] at (0.75,-0.3) {1};
    \draw[<->] (3.35,0) -- (3.35,1); \node[right] at (3.35,0.5) {1};
    \draw[<->] (3.35,1) -- (3.35,2); \node[right] at (3.35,1.5) {1};
    \node[below] at (1.5,-0.95) {俯视图};
  \end{scope}
\end{tikzpicture}
\end{center}
\nopagebreak[4]
（A）$\dfrac{11}2$；（B）$5$；（C）$\dfrac92$；（D）$4$。

\Needspace{4\baselineskip}
\textbf{解析}
\Needspace{5\baselineskip}
由三视图，该几何体是底面为俯视图所示六边形、高为 $1$ 的直棱柱。该六边形可看作长 $3$、宽 $2$ 的矩形去掉左右两个直角三角形（直角边 $1$、$1$），面积
\nopagebreak[4]
\[S=3\times2-2\times\frac12=4,\]
故体积 $V=4\times1=4$。选 D。

\gkanswer{答案：D。}

\end{answer}
\begin{answer}
\gksolhead{8}{}

\textbf{题目}：

椭圆 $\dfrac{x^2}{a^2}+\dfrac{y^2}{b^2}=1$（$a>b>0$）的左、右顶点分别是 $A$，$B$，左、右焦点分别是 $F_1$，$F_2$。若 $|AF_1|$，$|F_1F_2|$，$|F_1B|$ 成等比数列，则此椭圆的离心率为\mbox{（　　）}

（A）$\dfrac14$；（B）$\dfrac{\sqrt5}5$；（C）$\dfrac12$；（D）$\sqrt5-2$。

\Needspace{4\baselineskip}
\textbf{解析}
\Needspace{5\baselineskip}
$|AF_1|=a-c$，$|F_1F_2|=2c$，$|F_1B|=a+c$。由等比中项
\nopagebreak[4]
\[(2c)^2=(a-c)(a+c)=a^2-c^2,\]
即 $5c^2=a^2$，故 $e=\dfrac ca=\dfrac{\sqrt5}5$。选 B。

\gkanswer{答案：B。}

\end{answer}
\begin{answer}
\gksolhead{9}{}

\textbf{题目}：

已知 $f(x)=\sin^2\left(x+\dfrac\pi4\right)$。若 $a=f(\lg5)$，$b=f\left(\lg\dfrac15\right)$，则\mbox{（　　）}

（A）$a+b=0$；（B）$a-b=0$；（C）$a+b=1$；（D）$a-b=1$。

\Needspace{4\baselineskip}
\textbf{解析}
\nopagebreak[4]
\[f(x)=\sin^2\left(x+\frac\pi4\right)=\frac{1+\sin2x}2.\]
设 $t=\lg5$，则 $a=\dfrac{1+\sin2t}2$，$b=\dfrac{1+\sin(-2t)}2=\dfrac{1-\sin2t}2$，故 $a+b=1$。选 C。

\gkanswer{答案：C。}

\end{answer}
\begin{answer}
\gksolhead{10}{}

\textbf{题目}：

如图，$|OA|=2$（单位：$\mathrm m$），$|OB|=1$（单位：$\mathrm m$），$OA$ 与 $OB$ 的夹角为 $\dfrac\pi6$，以 $A$ 为圆心，$AB$ 为半径作圆弧 $\overset{\frown}{BDC}$ 与线段 $OA$ 的延长线交于点 $C$。甲、乙两质点同时从点 $O$ 出发，甲先以速率 $1$（单位：$\mathrm m/\mathrm s$）沿线段 $OB$ 行至点 $B$，再以速率 $3$（单位：$\mathrm m/\mathrm s$）沿圆弧 $\overset{\frown}{BDC}$ 行至点 $C$ 后停止；乙以速率 $2$（单位：$\mathrm m/\mathrm s$）沿线段 $OA$ 行至点 $A$ 后停止。设 $t$ 时刻甲、乙所到达的两点连线与它们经过的路径所围成图形的面积为 $S(t)$（$S(0)=0$），则函数 $y=S(t)$ 的图象大致是\mbox{（　　）}
\nopagebreak[4]
\begin{center}% 2012 年江西卷文科第 10 题：几何图形与 S(t) 的四个图象选项。
% 来源 PDF：普通高考/2012/2012江西文.pdf，第 1 页（印面 749）右栏。
% 几何：O(0,0)，A(2,0)，B 在 30° 方向且 OB=1，圆弧以 A 为圆心、AB≈1.24 为半径交 OA 延长线于 C。
% 正确图象 (B)：0≤t≤1 为抛物线，之后为直线（斜率更大），到 C 后水平。
\begin{tikzpicture}[scale=0.8, semithick, >=Stealth, font=\small]
  % ---------- 几何图形 ----------
  \begin{scope}
    \draw (0,0) node[below left] {$O$} -- (2,0) node[below] {$A$};
    \draw (0.866,0.5) node[above left] {$B$} -- (2,0);
    \draw (0,0) -- (0.866,0.5);
    \draw[thick] (0.866,0.5) arc[start angle=156.2, end angle=0, radius=1.239];
    \draw (2,0) -- (3.24,0) node[below] {$C$};
    \node[above] at (2.0,1.4) {$D$};
    \fill (0.866,0.5) circle (1.5pt);
    \fill (2,0) circle (1.5pt);
    \node[below] at (1.6,-1.35) {};
  \end{scope}
  % ---------- (A) ----------
  \begin{scope}[shift={(0.2,-3.9)}]
    \draw[->] (0,0) -- (3.2,0);
    \draw[->] (0,0) -- (0,2.4);
    \draw[smooth] (0,0) .. controls (0.6,0.05) and (1.2,0.9) .. (1.9,2.0) .. controls (2.2,2.25) and (2.5,2.3) .. (3.0,2.3);
    \node[below left] at (0,0) {$O$};
    \node at (1.6,-0.55) {(A)};
  \end{scope}
  % ---------- (B) ----------
  \begin{scope}[shift={(4.3,-3.9)}]
    \draw[->] (0,0) -- (3.2,0);
    \draw[->] (0,0) -- (0,2.4);
    \draw[smooth, domain=0:1] plot (\x, {0.5*\x*\x});
    \draw (1,0.5) -- (2.2,2.05) -- (3.0,2.05);
    \node[below left] at (0,0) {$O$};
    \node at (1.6,-0.55) {(B)};
  \end{scope}
  % ---------- (C) ----------
  \begin{scope}[shift={(0.2,-7.1)}]
    \draw[->] (0,0) -- (3.2,0);
    \draw[->] (0,0) -- (0,2.4);
    \draw[smooth] (0.3,0.05) .. controls (0.9,0.1) and (1.2,1.9) .. (1.9,2.25) .. controls (2.3,2.35) and (2.6,2.3) .. (3.0,2.3);
    \node[below left] at (0,0) {$O$};
    \node at (1.6,-0.55) {(C)};
  \end{scope}
  % ---------- (D) ----------
  \begin{scope}[shift={(4.3,-7.1)}]
    \draw[->] (0,0) -- (3.2,0);
    \draw[->] (0,0) -- (0,2.4);
    \draw (0.9,0) -- (0.9,2.2) -- (2.2,2.2);
    \node[below left] at (0,0) {$O$};
    \node at (1.6,-0.55) {(D)};
  \end{scope}
\end{tikzpicture}
\end{center}
\nopagebreak[4]
\Needspace{4\baselineskip}
\textbf{解析}
\Needspace{5\baselineskip}
当 $t\in[0,1]$ 时，甲在 $OB$ 上距 $O$ 为 $t$，乙在 $OA$ 上距 $O$ 为 $2t$，所围三角形面积
\nopagebreak[4]
\[S(t)=\frac12\cdot t\cdot2t\cdot\sin\frac\pi6=\frac12t^2,\]
图象是开口向上的抛物线。

\Needspace{5\baselineskip}
当 $t>1$ 时，乙已停在 $A$，甲沿以 $A$ 为圆心、$AB$ 为半径的圆弧运动，且运动方向与 $AP$ 垂直，故
\nopagebreak[4]
\[S'(t)=\frac12|AP|\cdot3=\frac32|AB|\]
为常数，图象是一段直线，且在 $t=1$ 处斜率由 $1$ 跳到 $\dfrac32|AB|>1$。甲到 $C$ 后 $S$ 为常数。

故图象为“抛物线—直线—水平线”，选 B。

\gkanswer{答案：B。}

\end{answer}
\begin{answer}
\gksolhead{11}{}

\textbf{题目}：

不等式 $\dfrac{x^2-9}{x-2}>0$ 的解集是\gkblank{}。

\Needspace{4\baselineskip}
\textbf{解析}
\Needspace{5\baselineskip}
$\dfrac{(x-3)(x+3)}{x-2}>0$，由序轴标根法得
\nopagebreak[4]
\[x\in(-3,2)\cup(3,+\infty).\]

\gkanswer{答案：$(-3,2)\cup(3,+\infty)$。}

\end{answer}
\begin{answer}
\gksolhead{12}{}

\textbf{题目}：

设单位向量 $\boldsymbol m=(x,y)$，$\boldsymbol b=(2,-1)$。若 $\boldsymbol m\perp\boldsymbol b$，则 $|x+2y|=$\gkblank{}。

\Needspace{4\baselineskip}
\textbf{解析}
\Needspace{5\baselineskip}
由 $\boldsymbol m\perp\boldsymbol b$ 得 $2x-y=0$，即 $y=2x$。又 $x^2+y^2=1$，故 $5x^2=1$，$|x|=\dfrac1{\sqrt5}$。于是
\nopagebreak[4]
\[|x+2y|=|5x|=\sqrt5.\]

\gkanswer{答案：$\sqrt5$。}

\end{answer}
\begin{answer}
\gksolhead{13}{}

\textbf{题目}：

等比数列 $\{a_n\}$ 的前 $n$ 项和为 $S_n$，公比不为 $1$。若 $a_1=1$，且对任意的 $n\in\mathbf N^*$，都有 $a_{n+2}+a_{n+1}-2a_n=0$，则 $S_5=$\gkblank{}。

\Needspace{4\baselineskip}
\textbf{解析}
\Needspace{5\baselineskip}
由 $a_{n+2}+a_{n+1}-2a_n=0$ 得 $q^2+q-2=0$，故 $q=1$ 或 $q=-2$。由 $q\neq1$ 得 $q=-2$，于是
\nopagebreak[4]
\[S_5=\frac{1-(-2)^5}{1-(-2)}=\frac{33}3=11.\]

\gkanswer{答案：$11$。}

\end{answer}
\begin{answer}
\gksolhead{14}{}

\textbf{题目}：

过直线 $x+y-2\sqrt2=0$ 上的点 $P$ 作圆 $x^2+y^2=1$ 的两条切线，若两条切线的夹角是 $60^\circ$，则点 $P$ 的坐标是\gkblank{}。

\Needspace{4\baselineskip}
\textbf{解析}
\Needspace{5\baselineskip}
设圆心为 $O$，两切线夹角为 $60^\circ$，则 $OP$ 平分该角，$\angle OPA=30^\circ$，故
\nopagebreak[4]
\[|OP|=\frac r{\sin30^\circ}=2.\]
\Needspace{5\baselineskip}
点 $P$ 在直线 $x+y-2\sqrt2=0$ 上且 $|OP|=2$，即 $x+y=2\sqrt2$，$x^2+y^2=4$。由 $(x+y)^2=8$ 得 $2xy=4$，$xy=2$，故 $x=y=\sqrt2$，即
\nopagebreak[4]
\[P(\sqrt2,\sqrt2).\]

\gkanswer{答案：$(\sqrt2,\sqrt2)$。}

\end{answer}
\begin{answer}
\gksolhead{15}{}

\textbf{题目}：

如图为某算法的程序框图，则程序运行后输出的结果是\gkblank{}。
\nopagebreak[4]
\begin{center}% 2012 年江西卷文科第 15 题：输出 T 的程序框图。
% 来源 PDF：普通高考/2012/2012江西文.pdf，第 1 页（印面 749）右栏。
% 流程：T=0,k=1 → sin(kπ/2)>sin((k-1)π/2)? 是 a=1，否 a=0 → T=T+a → k=k+1 → k<6? 是回到判断。
\begin{tikzpicture}[semithick, >=Stealth,
  block/.style={draw, minimum width=3.0cm, minimum height=0.62cm, align=center},
  dec/.style={draw, diamond, aspect=2.4, inner sep=1pt, align=center},
  term/.style={draw, rounded corners=6pt, minimum width=1.5cm, minimum height=0.52cm},
  io/.style={draw, trapezium, trapezium left angle=70, trapezium right angle=110,
             minimum width=2.2cm, minimum height=0.58cm},
  arr/.style={->}]
  \node (start) [term] at (0,8.6) {开始};
  \node (init)  [block] at (0,7.7) {$T=0,\ k=1$};
  \node (sin)   [block] at (0,6.2) {$\sin\dfrac{k\pi}2>\sin\dfrac{(k-1)\pi}2$?};
  \node (yes)   [block, right=1.5cm of sin] {$a=1$};
  \node (no)    [block, left=1.5cm of sin] {$a=0$};
  \node (add)   [block, below=3.2cm of sin] {$T=T+a$};
  \node (inc)   [block, below=0.55cm of add] {$k=k+1$};
  \node (dec2)  [dec, below=0.7cm of inc] {$k<6$?};
  \node (out)   [io, below=0.75cm of dec2] {输出 $T$};
  \node (end)   [term, below=0.5cm of out] {结束};
  \draw[arr] (start) -- (init);
  \draw[arr] (init) -- (sin);
  \draw[arr] (sin) -- node[above] {是} (yes);
  \draw[arr] (sin) -- node[above] {否} (no);
  \draw[arr] (yes.south) |- ($(add.east)+(0.0,0)$);
  \draw[arr] (no.south) |- ($(add.west)+(0.0,0)$);
  \draw[arr] (add) -- (inc);
  \draw[arr] (inc) -- (dec2);
  \draw[arr] (dec2) -- node[right] {否} (out);
  \draw[arr] (dec2.east) -- node[above] {是} ++(2.4,0) |- (sin.east);
  \draw[arr] (out) -- (end);
\end{tikzpicture}
\end{center}
\nopagebreak[4]
\Needspace{4\baselineskip}
\textbf{解析}
\Needspace{5\baselineskip}
$T=0$，$k=1$，当 $k<6$ 时比较 $\sin\dfrac{k\pi}2$ 与 $\sin\dfrac{(k-1)\pi}2$：
\nopagebreak[4]
\[k=1:\ 1>0\ \checkmark;\quad k=2:\ 0>1\ \times;\quad k=3:\ -1>0\ \times;\]
\nopagebreak[4]
\[k=4:\ 0>-1\ \checkmark;\quad k=5:\ 1>0\ \checkmark.\]
故 $T=1+0+0+1+1=3$。

\gkanswer{答案：$3$。}

\end{answer}
\begin{answer}
\gksolhead{16}{}

\textbf{题目}：

在 $\triangle ABC$ 中，角 $A$，$B$，$C$ 的对边分别为 $a$，$b$，$c$。已知 $3\cos(B-C)-1=6\cos B\cos C$。

（1） 求 $\cos A$；

（2） 若 $a=3$，$\triangle ABC$ 的面积为 $2\sqrt2$，求 $b$，$c$。

\Needspace{4\baselineskip}
\textbf{解析}
\Needspace{5\baselineskip}
\textbf{（1）} 由 $3\cos(B-C)-1=6\cos B\cos C$ 得
\nopagebreak[4]
\[3(\cos B\cos C+\sin B\sin C)-1=6\cos B\cos C,\]
\Needspace{5\baselineskip}
即 $3\sin B\sin C-3\cos B\cos C=1$，也就是
\nopagebreak[4]
\[-3\cos(B+C)=1.\]
因 $\cos(B+C)=-\cos A$，故 $\cos A=\dfrac13$。

\Needspace{5\baselineskip}
\textbf{（2）} $\sin A=\dfrac{2\sqrt2}3$。由面积
\nopagebreak[4]
\[S=\frac12bc\sin A=2\sqrt2\]
\Needspace{5\baselineskip}
得 $bc=6$。又由余弦定理
\nopagebreak[4]
\[9=b^2+c^2-2bc\cdot\frac13=b^2+c^2-4,\]
\Needspace{5\baselineskip}
故 $b^2+c^2=13$，$(b+c)^2=25$，$b+c=5$。于是 $b$，$c$ 是方程 $t^2-5t+6=0$ 的两根，即
\nopagebreak[4]
\[(b,c)=(2,3)\ \text{\text{或}}\ (3,2).\]

\gkanswer{答案（1）：$\dfrac13$。}

\gkanswer{答案（2）：$b=2$，$c=3$ 或 $b=3$，$c=2$。}

\end{answer}
\begin{answer}
\gksolhead{17}{}

\textbf{题目}：

已知数列 $\{a_n\}$ 的前 $n$ 项和 $S_n=kc^n-k$（其中 $c$，$k$ 为常数），且 $a_2=4$，$a_6=8a_3$。

（1） 求 $a_n$；

（2） 求数列 $\{na_n\}$ 的前 $n$ 项和 $T_n$。

\Needspace{4\baselineskip}
\textbf{解析}
\Needspace{5\baselineskip}
\textbf{（1）} 当 $n\geqslant2$ 时
\nopagebreak[4]
\[a_n=S_n-S_{n-1}=kc^n-kc^{n-1}=k(c-1)c^{n-1}.\]
\Needspace{5\baselineskip}
又 $a_1=S_1=k(c-1)$，与上式在 $n=1$ 时一致。由 $a_2=4$ 得 $k(c-1)c=4$；由 $a_6=8a_3$ 得 $c^3=8$，故 $c=2$，$k=2$。于是
\nopagebreak[4]
\[a_n=2^n.\]

\Needspace{5\baselineskip}
\textbf{（2）}
\nopagebreak[4]
\[T_n=\sum_{k=1}^nk\cdot2^k=(n-1)2^{n+1}+2.\]

\gkanswer{答案（1）：$a_n=2^n$。}

\gkanswer{答案（2）：$T_n=(n-1)2^{n+1}+2$。}

\end{answer}
\begin{answer}
\gksolhead{18}{}

\textbf{题目}：

如图，从 $A_1(1,0,0)$，$A_2(2,0,0)$，$B_1(0,1,0)$，$B_2(0,2,0)$，$C_1(0,0,1)$，$C_2(0,0,2)$ 这 $6$ 个点中随机选取 $3$ 个点。

（1） 求这 $3$ 点与原点 $O$ 恰好是正三棱锥的四个顶点的概率；

（2） 求这 $3$ 点与原点 $O$ 共面的概率。
\nopagebreak[4]
\begin{center}% 2012 年江西卷文科第 18 题：空间直角坐标系中的六个点 A1,A2,B1,B2,C1,C2。
% 来源 PDF：普通高考/2012/2012江西文.pdf，第 2 页（印面 750）第 1 栏。
% 投影：x 向左下、y 向右、z 向上。
\begin{tikzpicture}[semithick, >=Stealth, font=\small]
  \coordinate (O) at (0,0);
  \draw[->] (0,0) -- (-2.0,-1.0) node[below] {$x$};
  \draw[->] (0,0) -- (2.6,0) node[right] {$y$};
  \draw[->] (0,0) -- (0,2.6) node[above] {$z$};
  \fill (-1.0,-0.5) circle (1.8pt) node[below left] {$A_1$};
  \fill (-2.0,-1.0) circle (1.8pt) node[left] {$A_2$};
  \fill (1.0,0) circle (1.8pt) node[below] {$B_1$};
  \fill (2.0,0) circle (1.8pt) node[below] {$B_2$};
  \fill (0,1.0) circle (1.8pt) node[left] {$C_1$};
  \fill (0,2.0) circle (1.8pt) node[left] {$C_2$};
  \fill (0,0) circle (1.8pt) node[below left] {$O$};
\end{tikzpicture}
\end{center}
\nopagebreak[4]
\Needspace{4\baselineskip}
\textbf{解析}
从 $6$ 个点中任取 $3$ 个，共有 $\mathrm C_6^3=20$ 种。

\Needspace{5\baselineskip}
\textbf{（1）} 以 $O$ 为顶点、底面为正三角形的三棱锥，底面三点须到 $O$ 等距且两两等距。距离 $O$ 为 $1$ 的三点 $A_1,B_1,C_1$ 两两距离 $\sqrt2$，构成正三角形，$O$ 为其中心，是正三棱锥；$A_2,B_2,C_2$ 同理。其余取法到 $O$ 的距离不全相等，不构成以 $O$ 为顶点的正三棱锥。故
\nopagebreak[4]
\[P=\frac2{20}=\frac1{10}.\]

\textbf{（2）} 三点与 $O$ 共面：若含同一条坐标轴上的两个点，则这三点与 $O$ 共面（该轴过 $O$）。选一条轴（$3$ 种）取其两点，再从其余 $4$ 点中任取一点，共 $3\times4=12$ 种。不含同轴两点时，三点分别位于三条轴上，它们确定的平面不过 $O$，不共面。

\Needspace{5\baselineskip}
故
\nopagebreak[4]
\[P=\frac{12}{20}=\frac35.\]

\gkanswer{答案（1）：$\dfrac1{10}$。}

\gkanswer{答案（2）：$\dfrac35$。}

\end{answer}
\begin{answer}
\gksolhead{19}{}

\textbf{题目}：

如图，在梯形 $ABCD$ 中，$AB\parallel CD$，$E$，$F$ 是线段 $AB$ 上的两点，且 $DE\perp AB$，$CF\perp AB$，$AB=12$，$AD=5$，$BC=4\sqrt2$，$DE=4$。现将 $\triangle ADE$，$\triangle CFB$ 分别沿 $DE$，$CF$ 折起，使 $A$，$B$ 两点重合于点 $G$，得到多面体 $CDEFG$。

（1） 求证：平面 $DEG\perp$ 平面 $CFG$；

（2） 求多面体 $CDEFG$ 的体积。
\nopagebreak[4]
\begin{center}% 2012 年江西卷文科第 19 题：梯形 ABCD（图 1）与折叠后的多面体 CDEFG（图 2）。
% 来源 PDF：普通高考/2012/2012江西文.pdf，第 2 页（印面 750）第 2 栏。
% 图 1：AB=12，AE=3，EF=5，FB=4，DE=CF=4，AD=5，BC=4√2，CD=EF=5。
% 图 2：A、B 折起重合于 G，G 在 EF 的垂直平面内（EG=3，FG=4，DG=5，CG=4√2）。
\begin{tikzpicture}[scale=0.62, semithick, font=\small]
  % ---------- 图 1 ----------
  \begin{scope}
    \coordinate (A) at (0,0);
    \coordinate (E) at (3,0);
    \coordinate (F) at (8,0);
    \coordinate (B) at (12,0);
    \coordinate (D) at (3,4);
    \coordinate (C) at (8,4);
    \draw (A) -- (B) -- (C) -- (D) -- cycle;
    \draw (D) -- (E);
    \draw (C) -- (F);
    \node[below left] at (A) {$A$};
    \node[below] at (E) {$E$};
    \node[below] at (F) {$F$};
    \node[below right] at (B) {$B$};
    \node[above left] at (D) {$D$};
    \node[above right] at (C) {$C$};
    \node[below] at (6,-0.8) {图 1};
  \end{scope}
  % ---------- 图 2 ----------
  \begin{scope}[shift={(13.5,0.6)}]
    \coordinate (E2) at (0,0);
    \coordinate (F2) at (5,0);
    \coordinate (D2) at (0,4);
    \coordinate (C2) at (5,4);
    \coordinate (G2) at (1.8,-0.5);
    \draw (E2) -- (F2) -- (C2) -- (D2) -- cycle;
    \draw (D2) -- (E2);
    \draw (C2) -- (F2);
    \draw (G2) -- (E2);
    \draw (G2) -- (F2);
    \draw (G2) -- (D2);
    \draw (G2) -- (C2);
    \node[left] at (E2) {$E$};
    \node[right] at (F2) {$F$};
    \node[left] at (D2) {$D$};
    \node[right] at (C2) {$C$};
    \node[below] at (G2) {$G$};
    \node[below] at (2.5,-2.0) {图 2};
  \end{scope}
\end{tikzpicture}
\end{center}
\nopagebreak[4]
\Needspace{4\baselineskip}
\textbf{解析}
\textbf{（1）} 折叠前 $DE\perp AB$，$CF\perp AB$，$AE=\sqrt{5^2-4^2}=3$，$BF=\sqrt{(4\sqrt2)^2-4^2}=4$，故 $GE=3$，$GF=4$，$DG=5$，$CG=4\sqrt2$。

\Needspace{5\baselineskip}
取 $E$ 为原点，$EF$ 为 $x$ 轴，$ED$ 为 $y$ 轴。则 $D(0,4,0)$，$C(5,4,0)$，$F(5,0,0)$。设 $G(x,y,z)$，由
\nopagebreak[4]
\[GE=3,\quad GF=4,\quad DG=5\]
得 $x=1.8$，$y=0$，$z=2.4$，即 $G(1.8,0,2.4)$。

\Needspace{5\baselineskip}
平面 $DEG$ 的法向量
\nopagebreak[4]
\[\boldsymbol n_1=\overrightarrow{ED}\times\overrightarrow{EG}=(0,4,0)\times(1.8,0,2.4)=(9.6,0,-7.2),\]
\Needspace{5\baselineskip}
平面 $CFG$ 的法向量
\nopagebreak[4]
\[\boldsymbol n_2=\overrightarrow{FC}\times\overrightarrow{FG}=(0,4,0)\times(-3.2,0,2.4)=(9.6,0,12.8).\]
由 $\boldsymbol n_1\cdot\boldsymbol n_2=92.16-92.16=0$ 得两平面垂直。

\Needspace{5\baselineskip}
\textbf{（2）} 底面 $CDEF$ 是矩形，面积 $5\times4=20$，$G$ 到底面的高为 $z_G=2.4$。故
\nopagebreak[4]
\[V=\frac13\times20\times2.4=16.\]

\gkanswer{答案（1）：证明见解析。}

\gkanswer{答案（2）：$16$。}

\end{answer}
\begin{answer}
\gksolhead{20}{}

\textbf{题目}：

已知三点 $O(0,0)$，$A(-2,1)$，$B(2,1)$，曲线 $C$ 上任意一点 $M(x,y)$ 满足 $\left|\overrightarrow{MA}+\overrightarrow{MB}\right|=\overrightarrow{OM}\cdot\left(\overrightarrow{OA}+\overrightarrow{OB}\right)+2$。

（1） 求曲线 $C$ 的方程；

（2） 点 $Q(x_0,y_0)$（$-2<x_0<2$）是曲线 $C$ 上的动点，曲线 $C$ 在点 $Q$ 处的切线为 $l$，点 $P$ 的坐标是 $(0,-1)$，$l$ 与 $PA$，$PB$ 分别交于点 $D$，$E$，求 $\triangle QAB$ 与 $\triangle PDE$ 的面积之比。

\Needspace{4\baselineskip}
\textbf{解析}
\Needspace{5\baselineskip}
\textbf{（1）} $\overrightarrow{MA}=(-2-x,1-y)$，$\overrightarrow{MB}=(2-x,1-y)$，故
\nopagebreak[4]
\[\left|\overrightarrow{MA}+\overrightarrow{MB}\right|=2\sqrt{x^2+(y-1)^2}.\]
\Needspace{5\baselineskip}
又 $\overrightarrow{OM}\cdot(\overrightarrow{OA}+\overrightarrow{OB})=(x,y)\cdot(0,2)=2y$，故
\nopagebreak[4]
\[2\sqrt{x^2+(y-1)^2}=2y+2,\]
\Needspace{5\baselineskip}
即 $x^2+(y-1)^2=(y+1)^2$，化简得
\nopagebreak[4]
\[x^2=4y.\]

\Needspace{5\baselineskip}
\textbf{（2）} 设 $Q(x_0,y_0)$，则 $y_0=\dfrac{x_0^2}4$。切线 $l$：
\nopagebreak[4]
\[y=\frac{x_0}2x-\frac{x_0^2}4.\]
直线 $PA:y=-x-1$，$PB:y=x-1$。

\Needspace{5\baselineskip}
由 $l$ 与 $PA$ 得 $x_D=\dfrac{x_0-2}2$，由 $l$ 与 $PB$ 得 $x_E=\dfrac{x_0+2}2$。于是
\nopagebreak[4]
\[S_{\triangle PDE}=\frac12\left|\begin{vmatrix}x_D&y_D+1\\x_E&y_E+1\end{vmatrix}\right|=\frac{4-x_0^2}4.\]
\Needspace{5\baselineskip}
而 $\triangle QAB$ 以 $AB=4$ 为底，高为 $\left|y_0-1\right|=\dfrac{4-x_0^2}4$，故
\nopagebreak[4]
\[S_{\triangle QAB}=\frac12\times4\times\frac{4-x_0^2}4=\frac{4-x_0^2}2.\]
所以面积之比为 $2$。

\gkanswer{答案（1）：$x^2=4y$。}

\gkanswer{答案（2）：$2$。}

\end{answer}
\begin{answer}
\gksolhead{21}{}

\textbf{题目}：

已知函数 $f(x)=(ax^2+bx+c)\mathrm e^x$ 在 $[0,1]$ 上单调递减且满足 $f(0)=1$，$f(1)=0$。

（1） 求 $a$ 的取值范围；

（2） 设 $g(x)=f(x)-f'(x)$，求 $g(x)$ 在 $[0,1]$ 上的最大值和最小值。

\Needspace{4\baselineskip}
\textbf{解析}
\Needspace{5\baselineskip}
\textbf{（1）} 由 $f(0)=1$ 得 $c=1$；由 $f(1)=0$ 得 $(a+b+1)\mathrm e=0$，即 $b=-1-a$。于是
\nopagebreak[4]
\[f'(x)=\mathrm e^x\left[ax^2+(a-1)x-a\right].\]
\Needspace{5\baselineskip}
设 $h(x)=ax^2+(a-1)x-a$，需 $h(x)\leqslant0$ 在 $[0,1]$ 上恒成立。由 $h(0)=-a\leqslant0$ 得 $a\geqslant0$，由 $h(1)=a-1\leqslant0$ 得 $a\leqslant1$。当 $0\leqslant a\leqslant1$ 时 $h$ 为凸函数，最大值在端点取得，故条件等价于
\nopagebreak[4]
\[a\in[0,1].\]

\Needspace{5\baselineskip}
\textbf{（2）}
\nopagebreak[4]
\[g(x)=f(x)-f'(x)=\left[1+a-2ax\right]\mathrm e^x.\]
当 $a=0$ 时 $g(x)=\mathrm e^x$，最大值 $\mathrm e$，最小值 $1$。

\Needspace{5\baselineskip}
当 $0<a\leqslant1$ 时，
\nopagebreak[4]
\[g'(x)=(1-a-2ax)\mathrm e^x.\]
若 $0<a<\dfrac13$，$g'>0$，$g$ 递增，最大值 $g(1)=(1-a)\mathrm e$，最小值 $g(0)=1+a$。

\Needspace{5\baselineskip}
若 $\dfrac13\leqslant a\leqslant1$，$g$ 在 $x_0=\dfrac{1-a}{2a}$ 处取最大值
\nopagebreak[4]
\[g(x_0)=2a\mathrm e^{\frac{1-a}{2a}},\]
最小值在端点取得，为 $\min\{1+a,(1-a)\mathrm e\}$，分界点为 $a=\dfrac{\mathrm e-1}{\mathrm e+1}$。

综上：最大值当 $0<a\leqslant1$ 时为 $2a\mathrm e^{\frac{1-a}{2a}}$（$a=0$ 时为 $\mathrm e$）；最小值当 $0\leqslant a\leqslant\dfrac{\mathrm e-1}{\mathrm e+1}$ 时为 $1+a$，当 $\dfrac{\mathrm e-1}{\mathrm e+1}\leqslant a\leqslant1$ 时为 $(1-a)\mathrm e$。

\gkanswer{答案（1）：$[0,1]$。}

\gkanswer{答案（2）：最大值 $2a\mathrm e^{\frac{1-a}{2a}}$（$0<a\leqslant1$），$a=0$ 时为 $\mathrm e$；最小值当 $0\leqslant a\leqslant\dfrac{\mathrm e-1}{\mathrm e+1}$ 时为 $1+a$，当 $\dfrac{\mathrm e-1}{\mathrm e+1}\leqslant a\leqslant1$ 时为 $(1-a)\mathrm e$。}

\end{answer}